\section{The full labelled arithmetic atlas}
Fix a prime $p=4m+1$. The complete first-denominator interval is
\[
 m+1\le a\le3m,\qquad R=4a-p.
\]
It has $3\le R\le2p-3$ and $\gcd(a,R)=\gcd(p,R)=1$.
For $a=\prod_q q^{e_q}$ retain every vector
$f\in\prod_q\{0,\ldots,2e_q\}$ and $u=\prod_q q^{f_q}$.
The inverse is $f_q=v_q(u)$.
The channel formulas are
\begin{align}
 Q_E&=\left(a,\frac{pa+a^2/u}{R},\frac{pa+p^2u}{R}\right),\\
 Q_M&=\left(a,\frac{p(a+a^2/u)}R,\frac{p(a+u)}R\right).
\end{align}
For all raw orders, use the disjoint labels $(E,u,0),(E,u,1),(M,u,0)$;
the exterior bit reverses the last two positions. The middle reversal is
already $u\mapsto a^2/u$, so another middle bit would duplicate it.
The first coordinate need not be the smallest at an arbitrary rational
candidate; it was the smallest only in the completeness argument below.

\begin{lemma}[Exact rational and integral return]\label{lem:denom}
Every such candidate is a positive rational decomposition of $4/p$.
Its last two entries have the same reduced denominator
\[
 D_E=\frac R{\gcd(R,4u+1)},\qquad
 D_M=\frac R{\gcd(R,u+a)}.
\]
In particular $D$ is odd and prime to $p$, and $D=1$ exactly at an integral
candidate. ES holds at $p$ exactly when some full-atlas candidate has $D=1$.
The number of full ordered labels is at most $\tfrac32(p-1)^2$.
\end{lemma}
\begin{proof}
Writing $Q=(a,Y,Z)$, substitution gives
$(RY-pa)(RZ-pa)=p^2a^2$; both factors are positive.
Expansion proves the rational reciprocal identity.
For $E$, multiplying the $Y$ numerator by the unit $u$ gives
$a^2(4u+1)$ modulo $R$, and the $Z$ numerator is
$4a^2(4u+1)$ modulo $R$. For $M$, the analogous numerators reduce, up to
units, to $u+a$. Their gcds with $R$ therefore agree exactly, proving the
minimal denominators.

For any positive integral solution, move a smallest denominator to the
first position and retain the permutation. Then $p/4<a\le3p/4$, so it is
in this interval. The positive factors $d_Y=RY-pa,d_Z=RZ-pa$ multiply to
$p^2a^2$. As $p\nmid a$, $v_p(d_Y)$ is 0,1 or 2. The three inverse labels
are respectively
\[
 (E,a^2/d_Y,0),\qquad (M,pa^2/d_Y,0),\qquad(E,d_Y/p^2,1).
\]
The congruence $d_Y\equiv-pa\pmod R$ gives the corresponding gate by
multiplication by units. Conversely an integral candidate is a solution.
Finally $\tau(a^2)\le2a-1$ by pairing divisors around $a$; hence
$3\sum_{a=m+1}^{3m}\tau(a^2)\le24m^2=\tfrac32(p-1)^2$.
This is the full-shell form of the inherited Type I/II divisor
parametrizations \cite[\S2]{et}; it is not claimed as a new parametrization.
\end{proof}

For every rational candidate keep
\[
 g=\gcd(a,u),\quad h=g^2/u,\quad r=u/g,\quad s=a/g,
\]
so $h,r,s$ are positive integers, $a=hrs$, $u=hr^2$ and $\gcd(r,s)=1$.
The quotients $\kappa=(pr+s)/R$ and $\lambda=(r+s)/R$ remain rational
until the gate passes. In the unswapped order the inverse is
$u=a^2/(RY-pa)$ for $E$ and $u=pa^2/(RY-pa)$ for $M$.
Clearing $X=DQ$ gives a positive integral solution at $pD$, not at $p$:
\[
 F_p(X)=4X_0X_1X_2-p\sum_{i<j}X_iX_j
       =p(D-1)\sum_{i<j}X_iX_j.
\]
The prime, the scale and the raw order therefore remain marks.

\section{Arithmetic obstruction groups become geometric cocycles}
For a candidate $\alpha$ let $g_E=4u+1$ or $g_M=u+a$, $h_0=\gcd(R,g_c)$,
and $D=R/h_0$. Its cyclic obstruction group is the \emph{actual} subgroup
\[
 O_\alpha=\langle[g_c]_R\rangle\le\Z/R.
\]
The map $\Z/D\to O_\alpha$, $k\mapsto kg_c\bmod R$, is a bijection.
On a representative $v\in O_\alpha\cap\{0,\ldots,R-1\}$ its inverse is
\[
 k=\frac v{h_0}(g_c/h_0)^{-1}\pmod D;
\]
modulo one the inverse is the unique zero class.
Each prime power is retained by
\[
 v_\ell(D)=\max\{v_\ell(R)-v_\ell(g_c),0\}.
\]
Combining this inverse with Theorem~\ref{thm:class} gives the \emph{faithful}
map
\begin{equation}\label{eq:cohomap}
 O_\alpha\xrightarrow{\sim}\langle[b_D]\rangle,\qquad
 kg_c\longmapsto k[b_D].
\end{equation}
The inverse first reads minus the third cusp evaluation modulo $D$, then
multiplies by $g_c$ modulo $R$. Thus the denominator obstruction is not
merely assigned a similarly sized group: it is explicitly identified
with a split cyclic subgroup of the geometric cocycle group.

\begin{definition}
The \emph{full denominator-labelled cover} at $p$ is the disjoint union,
over every ordered atlas label $\alpha$, of the covers with fibres
$\{\alpha\}\times\mathcal F_{D_\alpha}$ and actions
$(\alpha,v)\mapsto(\alpha,a_jv)$. Its decoration includes the prime,
shell, factorization, exponent vector, channel, swap, $h,r,s$, channel
quotient, rational triple, clearing scale and all local losses.
\end{definition}

\begin{theorem}[Complete failure object]\label{thm:failure}
ES fails at $p$ if and only if this full labelled cover has no section.
Equivalently, every summand has a nonzero cocycle class of exact order
$D_\alpha\ge3$, and every cusp permutation has no fixed sheet.
The compactified components, degrees, ramification and genera of the
complete object are given by Theorem~\ref{thm:orbits}, applied with their
full arithmetic multiplicities.
\end{theorem}
\begin{proof}
Over the connected base a section has a constant discrete candidate label.
Theorem~\ref{thm:class} says its value in that summand is single-valued
exactly when $D_\alpha=1$. Lemma~\ref{lem:denom} proves the arithmetic
biconditional. The finite connected components are precisely the orbits,
so Theorem~\ref{thm:orbits} supplies the stated complete classification.
\end{proof}

The cover degree is bounded by
$\tfrac32(p-1)^2(2p-3)^3$. The small input consisting of all labelled
$D_\alpha$ and the three formulas specifies every sheet and every edge;
it is not necessary to expand that many sheets to validate the atlas.
An actual global negative certificate would have to provide \emph{all}
labels and a nonzero prime-power loss on each. No global negative example
is supplied. The algorithm returns an existing positive witness on all
four primes used in its full-atlas tests.

The previous operator $H_p=\operatorname{diag}((D_\alpha-1)/2)$ has a
separate meaning. Its inverse, when all $D>1$, multiplies the $\alpha$-th
coordinate by $2/(D_\alpha-1)$. Every geometric loop permutation is
invertible regardless of whether $D=1$. Neither assertion is the inverse
of the other. The same denominator connects them, with its label retained.

\paragraph{Restricted prototype.}
At $p=241,a=64,R=15$ every divisor is $u=2^e$, $0\le e\le12$.
For each canonical channel, complete modular arithmetic gives
$D=3$ four times, $D=5$ six times, and $D=15$ three times.
Including the exterior swap gives 39 labels, with multiplicities
12,18,9 respectively. The associated cover has
\[
 32949\text{ sheets},\qquad60\text{ connected components},\qquad
 \sum g=5448.
\]
All 39 summands are sectionless. This is an actual complete negative
\emph{shell}, not a false prime: other shells at 241 succeed.
Taking the gcd of its denominators gives one and would destroy the
leafwise obstruction. The validator rejects a restricted certificate,
a deleted row, or a full atlas containing a successful candidate.


\paragraph{Existing arithmetic transports lift, but do not erase their labels.}
Exterior order reversal and middle complementation preserve $D$.
The retained relation from $E(u)$ to $M(u')$, where $u,u'\mid a^2$ and
$u'\equiv pu\pmod R$, preserves $D$ because
$u'+a\equiv a(4u+1)\pmod R$ and $a$ is a unit.
Thus every such labelled edge has the explicit equivariant cover map
$(\alpha,v)\mapsto(\alpha',v)$, with reverse edge as inverse.
The endpoint labels and the whole chosen path are retained; forgetting
them would make this cover-valued functor nonfaithful on paths.
The 39-state restricted graph above has three components, of sizes
12,18,9 at levels 3,5,15. The certificate lists all edges.
These particular transports are defect-preserving, not a claimed
strict descent. Cross-shell arithmetic still needs an additional
argument; arbitrary independent choices of $D_\alpha$ are not substituted
for the simultaneously calculated values from one prime.
